\documentclass{article}
\usepackage{amsmath}

\begin{document}

\section{Base model analysis}

The corrected base model is
\[
\begin{aligned}
S' &= S\left[r_S\left(1-\frac{S+\alpha R}{K}\right)
- \lambda(t)\right], \\
R' &= Rr_R\left(1-\frac{\beta S+R}{K}\right).
\end{aligned}
\]

Here, $S(t)$ is the sensitive cancer cell population and
$R(t)$ is the resistant cancer cell population at time $t$.
Both populations consist of cancer cells. The difference is
that the treatment affects sensitive cells, while resistant
cells are assumed to be unaffected by it directly.

The idea of adaptive therapy is that these two populations
compete for resources. Treatment reduces the sensitive
population, but this also frees space and resources for
resistant cells. This is called competitive release.
The goal is therefore to maintain enough sensitive cells
to limit resistant growth, while keeping the total cancer
population under control.

The equations are based on assumptions about cell growth
rather than derived directly from biological laws.
We first assume that, without treatment or resource limits,
the growth of a population is proportional to its size:
\[
S(t+\Delta t)-S(t)
\approx r_S S(t)\Delta t.
\]
Dividing by $\Delta t$ and taking the limit gives
\[
S'=r_SS.
\]
This gives exponential growth:
\[
S(t)=S(0)e^{r_St}.
\]
The parameter $r_S$ is the growth rate per cell under these
conditions. The same argument applies to $R$, with growth
rate $r_R$.

However, exponential growth assumes that resources do not
run out. To account for limited resources, we introduce
the carrying capacity $K$. For a single population, we write
\[
S'=r_SS\left(1-\frac{S}{K}\right).
\]
The idea is that growth slows down as the population gets
closer to $K$. If $0<S<K$, the population grows.
At $S=K$, its net growth is zero. If $S>K$, it decreases.

Next, we include competition between the two populations.
Sensitive cells do not only compete with other sensitive
cells, but also with resistant cells. We therefore replace
$S$ in the crowding term by
\[
S+\alpha R.
\]
Here, $\alpha$ describes how strongly resistant cells
affect the sensitive population compared to sensitive
cells themselves. For example, if $\alpha=2$, one resistant
cell contributes twice as much to this crowding term as
one sensitive cell.

Similarly, resistant cells experience crowding through
\[
\beta S+R.
\]
So $\beta$ describes the effect of sensitive cells on
resistant cells. Increasing $\beta$ makes sensitive cells
stronger competitors against resistant cells.
Increasing $\alpha$ has the opposite direction of influence:
it makes resistant cells stronger competitors against
sensitive cells.

Finally, treatment removes sensitive cells at rate
$\lambda(t)S$. We subtract this from their growth:
\[
S'=r_SS\left(1-\frac{S+\alpha R}{K}\right)-\lambda(t)S.
\]
There is no corresponding treatment term in the equation
for $R$, because complete resistance is assumed.

The parameters $r_S$, $r_R$, and $\lambda$ have units of
inverse time. The parameters $\alpha$ and $\beta$ are
dimensionless, while $K$ has the same units as $S$ and $R$.

\subsection{When do the populations grow or decrease?}

For $S>0$, we can divide the first equation by $S$:
\[
\frac{S'}{S}
=
r_S\left(1-\frac{S+\alpha R}{K}\right)-\lambda(t).
\]
This gives the relative growth rate of the sensitive
population. The population decreases precisely when
\[
\lambda(t)>
r_S\left(1-\frac{S+\alpha R}{K}\right).
\]
So treatment being present does not automatically mean
that the sensitive population decreases. The treatment
loss must be larger than its growth.

For $R>0$ and $r_R>0$, the second equation gives
\[
\begin{aligned}
R'>0 &\iff \beta S+R<K,\\
R'=0 &\iff \beta S+R=K,\\
R'<0 &\iff \beta S+R>K.
\end{aligned}
\]
This shows how sensitive cells can limit resistant growth.
A larger sensitive population increases the crowding
experienced by resistant cells. If this crowding is strong
enough, the resistant population decreases.

For the baseline choice $\alpha=\beta=1$, both populations
experience the same crowding factor:
\[
1-\frac{S+R}{K}.
\]
In particular, if $S+R<K$, the resistant population still
grows. Sensitive cells can slow this growth down, but do
not make it negative in this region.

There is also a useful boundary observation. At $S+R=K$,
\[
(S+R)'=-\lambda(t)S\leq 0
\]
when $\lambda(t)\geq 0$. Therefore, a solution starting
with $S+R\leq K$ cannot cross this boundary upwards.
For this baseline case, the resistant population cannot
decrease along such a solution. This limits what the
base model can achieve: it can describe delaying resistant
growth without describing its elimination.

\subsection{Adaptive treatment}

To model adaptive treatment, $\lambda(t)$ changes between
zero and a fixed treatment rate $\lambda_{\mathrm{on}}$.
Let
\[
N(t)=S(t)+R(t), \qquad N_0=N(0).
\]

In our simulation, treatment starts on. It is switched off
when $N(t)$ falls to half its initial value and restarted
when $N(t)$ returns to its initial value:
\[
\begin{aligned}
N(t)\leq \tfrac12N_0
&\quad\Rightarrow\quad \text{switch treatment off},\\
N(t)\geq N_0
&\quad\Rightarrow\quad \text{switch treatment on}.
\end{aligned}
\]
Between these thresholds, the previous treatment state
is kept.

The idea is to avoid removing too many sensitive cells.
When treatment stops, sensitive cells can recover and
increase competition for resistant cells. Whether this
controls the total population for longer depends on the
parameters and initial conditions, so we compare it with
continuous treatment and no treatment.

\subsection{Assumptions and limitations}

The model describes population sizes, not the positions
or individual behaviour of cells. Resource competition
is included through the crowding terms, but the resources
themselves are not modelled.

There is also no term describing sensitive cells becoming
resistant. If $R(0)=0$, then $R(t)=0$ for all later times.
We therefore need $R(0)>0$ to study resistance in this
model. What we study is the growth of an existing resistant
population, rather than the appearance of new resistance.

\end{document}